% ============================================================
%  封面信息
% ============================================================
\school{计算机科学与技术学院}
\major{软件工程（软件工作）}
\student{2251321}{徐嘉豪}
\thesistitle{一款轻量级实时数字协作白板系统}{设计与开发}
\thesistitleeng{A Lightweight Real-time Digital Collaborative Whiteboard System}{Design and Development}
\thesisadvisor{张晶\quad 指导教师}
\thesisdate{\the\year}{\the\month}{\the\day} % 使用当前日期; 也可以自定义日期

% ============================================================
%  信息说明页
% ============================================================

% 成果类型：design（毕业设计）或 thesis（毕业论文）
\infotype{thesis}

% 内容简述（请用300字以内简要概述）
\infoabstract{%
QuickBoard 是一个“打开即用”的在线协作白板系统，面向课堂讨论与远程协作场景，目标是在无需注册的前提下，实现秒级创建/加入房间与稳定的多人实时绘制体验。系统前端基于 Vue3 与 Fabric 实现可编辑画布，提供选择、画笔、形状、橡皮与导出等能力；后端基于 Node.js 与 Socket.IO 完成房间管理与实时消息通信，并采用服务端权威裁决与增量同步策略减少冲突与回滚；同时接入独立 OCR 服务识别数学公式并在前端以 KaTeX 渲染与可编辑流程实现“识别—确认—插入—编辑”闭环。本文给出需求分析、总体架构与关键实现，并通过单元测试与构建验证及场景化功能测试评估系统可用性与稳定性。
}

% 毕业设计：图纸数量与正文字数（成果类型为 design 时填写，两者须同时填写）
% \infodrawings{5}
% \infowordcount{15000}

% 毕业论文：正文总字数（成果类型为 thesis 时填写）
\infothesiswords{8000}

% 随附资料（每条用 \item 开头；不填则显示空白行）
\infomaterials{
  \item 程序源代码（前端、后端与 OCR 服务）；
  \item 测试报告与构建报告（\texttt{doc/test\_reports}）。
}
